\begin{proof}[Proof of Corollary~\ref{thm:stronglaw} for the other distributions]
Multiplying every coefficient by a nonzero constant does not change the
zeros of $f_n$, and a nondegenerate centrally symmetric distribution has
positive second moment, since it is not concentrated at zero. We may
therefore assume $\mathbb E|\xi_0|^2=1$ for a bounded distribution, as in
Proposition~\ref{prop:coefficient-transfer}, and this holds for the other
three distributions.
We follow the proof for Rademacher coefficients in
Section~\ref{sec:unit-circle} step by step, with the estimates of
Proposition~\ref{prop:coefficient-transfer} instead of those of
Sections~\ref{sec:concentration} and~\ref{sec:local-roots}, and we write
out every step in which a constant, a degree or an exceptional
probability changes.
By Proposition~\ref{prop:coefficient-transfer}, the occupation,
logarithmic, energy, annular mesh, local-count and maximal-tail estimates
hold, the logarithmic estimate on the energy event of
that proposition when the coefficient distribution requires it, and the
additional exceptional probabilities are
exponential in $N$ or explicitly summable.
Since $\mathbb E|\xi_0|^2=1$, the normalization
$\sigma_N(r)^2=\sum_{k=0}^Nr^{2k}$ is that of the sign case, so that, as
in Proposition~\ref{prop:profiles}, the variance profile is
\[
 F(x)=\frac12\log\int_0^1e^{2xt}\dd t,\qquad F'(0)=\frac12,
\]
and Lemma~\ref{lem:profile-rate} applies unchanged.
Together with the matching argument of Section~\ref{sec:unit-circle},
these estimates yield the conclusion for nonzero coefficients.
In what follows, let $\Pi_K(N)$ be the right-hand side of the
concentration bound of Proposition~\ref{prop:coefficient-transfer}
without its term $\Prob(E^c)$, with the constants of the distribution, and
allow $C(K)$ to depend on the distribution.

\emph{Step 1.}
Proposition~\ref{prop:coefficient-transfer} implies that
Lemmas~\ref{lem:second-derivative} and~\ref{lem:tails} hold for Steinhaus
and Gaussian coefficients with the same constants, for Gaussian
coefficients on the envelope event $\max_{k\le2N}|\xi_k|\le N$.
Hence $\rho$, $a$, $M_2$ and $d_0$ and the bounds
\eqref{eq:unit-block-length}--\eqref{eq:unit-parameters} are unchanged.
For bounded coefficients, the proof of
Proposition~\ref{prop:coefficient-transfer} yields the tail amplitude
$Ba$, with $K+2$ in place of $K$ and $m_{\max}=m_j$, and the
second-derivative bound $BM_2$.
Since \eqref{eq:unit-block-length} does not involve the coefficients,
\eqref{eq:unit-amplitude}--\eqref{eq:unit-parameters} become
\begin{align*}
 Ba&\le45Be^{2(K+2)}N^{7/16}\sqrt{\log N},\\
 \frac{2Ba}{d_0}&\le90Be^{2(K+2)}N^{-67/64}\sqrt{\log N},\\
 \frac{2NBa}{d_0}&\le90Be^{2(K+2)}N^{-3/64}\sqrt{\log N},\\
 \frac{4(BM_2)(Ba)}{d_0^2}
 &\le180\cdot100B^2e^{K+4}e^{2(K+2)}N^{-1/32}\log N,
\end{align*}
which are those bounds multiplied by $B$, $B$, $B$ and $B^2$.
The right-hand sides of the last two lines tend to zero, so that
$2Ba/d_0\le1/N$ and $4(BM_2)(Ba)\le d_0^2$ for all $N$ large enough,
depending on $K$ and $B$.

\emph{Step 2.}
For Steinhaus and Gaussian coefficients this step is the same as before,
since $a$, $M_2$ and $d_0$ are unchanged.
For bounded coefficients, we apply Lemma~\ref{lem:isolation-interface} to
each good annular root $\alpha$, with $P=f_N$, the amplitude $Ba$,
$|f_N'(\alpha)|\ge d_0$ and $|f_N''|\le BM_2$ on the disk of radius
$2Ba/d_0$ about $\alpha$.
This disk lies in $|z|\le1+(K+1)/N$, since $2Ba/d_0\le1/N$, and we have
$4(BM_2)(Ba)\le d_0^2$ by Step~1.
The lemma then shows that the disk contains exactly one zero and that
$|f_N|\ge3Ba/2$ on its boundary, and we choose the common regular level in
$(Ba,3Ba/2)$.

\emph{Step 3.}
The proof of Corollary~\ref{cor:radial-probability} relies on
the deterministic Lemmas~\ref{lem:annular-mass} and~\ref{lem:thin-band},
which require a nonzero polynomial of degree at most $N$.
For Steinhaus coefficients $f_N(0)=\xi_0\ne0$, and for Gaussian
coefficients $f_N(0)=\xi_0\ne0$ almost surely.
For bounded coefficients, $f_N\equiv0$ implies $V=0$ at every radius, so
that the zero polynomial lies in the energy exception $E^c$.
It follows that, with the concentration bound of
Proposition~\ref{prop:coefficient-transfer} instead of
Theorem~\ref{thm:sign-logarithms}, the corollary holds for each distribution, with
its constants, outside four events at the width $K$ and five events at the
width $2$, of probabilities at most $\Pi_K(N)+\Prob(E^c)$ and
$\Pi_2(N)+\Prob(E^c)$, with $\Prob(E^c)$ at the widths $K$ and $2$.
Using the bound \eqref{eq:gauss-small-derivatives} or
\eqref{eq:bounded-small-derivatives} on $\Prob\{B_{K,N}>N^{31/32}\}$ from
the proof of Proposition~\ref{prop:coefficient-transfer} instead of
Lemma~\ref{lem:small-derivatives} (for Steinhaus coefficients
\eqref{eq:annular-small-derivatives} still holds), we obtain
\eqref{eq:unit-unselected} with the constant $C(K)$ of the distribution.
For bounded coefficients, the crossing bound \eqref{eq:unit-crossings}
needs $2Ba/d_0<N^{-65/64}$ instead of $2a/d_0<N^{-65/64}$, and by Step~1
we have
\[
 N^{65/64}\,\frac{2Ba}{d_0}
 \le90Be^{2(K+2)}N^{-67/64+1+1/64}\sqrt{\log N}
 =90Be^{2(K+2)}N^{-1/32}\sqrt{\log N}.
\]
This tends to zero, as \eqref{eq:unit-radius-ratio} does.
The bound on $c_N/N$ and the case $q=0$ of \eqref{eq:unit-thin-band}
then hold with the constant $C(2)$ of the distribution.

\emph{Step 4.}
As in Step~4 of Section~\ref{sec:unit-circle}, outside one maximal tail
event we have $|f_n-f_N|<a$, or $|f_n-f_N|<Ba$ for bounded coefficients,
on every selected boundary.
Theorem~\ref{thm:T4} needs nonzero polynomials, and outside the events of
Step~3 we have $f_N\not\equiv0$, which implies
$f_n\not\equiv0$, since its first $N+1$ coefficients are those of $f_N$.
If $\Prob(\xi_0=0)=0$, every coefficient is nonzero almost surely, so
that $\deg f_n=n$ and $\deg f_N=N$, and Step~4 holds as before, with
\eqref{eq:unit-block}.

Suppose now that $\Prob(\xi_0=0)>0$.
Then the zero-polynomial event is contained in the amplitude exception.
Since the probabilities of the trailing-run events below are summable
along $N=j^8$ once $8D|\log\Prob(\xi_0=0)|>1$, any
$D>1/(8|\log\Prob(\xi_0=0)|)$ would do.
We choose $D>4/|\log\Prob(\xi_0=0)|$, which keeps a margin.
Let the trailing-run event be the event that the trailing run of zero
coefficients at $N=j^8$ exceeds $\lceil D\log N\rceil$.
Since $0<\Prob(\xi_0=0)<1$, this event has probability at most
$\Prob(\xi_0=0)^{\lceil D\log N\rceil}\le\Prob(\xi_0=0)^{D\log N}
 =N^{-D|\log\Prob(\xi_0=0)|}$, and, since $8D|\log\Prob(\xi_0=0)|>32$,
we also have $\sum_{j=1}^\infty j^{-8D|\log\Prob(\xi_0=0)|}<\infty$.
Outside the trailing-run event we have
$\deg f_N\ge N-\lceil D\log N\rceil$ and $\deg f_n\le n$,
so that, for every $N\le n\le N+m_j$,
\begin{align*}
 \deg f_n-\deg f_N&\le n-N+\lceil D\log N\rceil,\\
 0\le\frac{\deg f_n-\deg f_N}{N}
 &\le\frac{m_j}{N}+\frac{\lceil D\log N\rceil}{N}
 \longrightarrow0,
\end{align*}
where $0\le$ holds since the prefix degrees are nondecreasing, and the
limit follows from $n-N\le m_j$ and \eqref{eq:unit-block-length}.
Thus the actual-degree term in Theorem~\ref{thm:T4}, divided by $N$,
tends to zero.
Since each selected domain contains exactly one root of $f_N$ and $U_N$
counts the other roots of $f_N$, we still have $\deg f_N-s=U_N$, although
$\deg f_N$ may now be smaller than $N$.
Hence, by Theorem~\ref{thm:T4} and the degree bound above, for
$N\le n\le N+m_j$ we have
\begin{align*}
 |\nu_n(1)-\nu_N(1)|
 &\le c_N+(\deg f_N-s)+\max(0,\deg f_n-\deg f_N)\\
 &\le c_N+U_N+m_j+\lceil D\log N\rceil,\\
 \left|\frac{\nu_n(1)}n-\frac12\right|
 &\le\left|\frac{\nu_N(1)}N-\frac12\right|
       +\frac{c_N+U_N}{N}+\frac{3m_j}{2N}
       +\frac{\lceil D\log N\rceil}N,
\end{align*}
where the third line follows from the second as in the chain of
\eqref{eq:unit-block}, with the second line in place of the matching bound
there and $\lceil D\log N\rceil/n\le\lceil D\log N\rceil/N$.
The initial zero multiplicity, the index of the first nonzero
coefficient, is finite almost surely, since the coefficients are
independent and $\Prob(\xi_0\ne0)>0$, and so, divided by $N$, it tends to
zero, and the bound on $U_N/N$ already counts these zeros at the origin,
since they are unselected roots.
There are only finitely many zero polynomials, and their counts do not
affect the limit.

\emph{Step 5.}
We now also include in the event $E_{K,j}$ the energy exceptions at the
four radii of the width $K$ and at the five radii of the width $2$, the
additional local-count exception and, for Gaussian coefficients, the
complement of the envelope event, which the derivative and tail bounds
share.
For bounded distributions with an atom at zero we also include the trailing-run
event.
For Steinhaus coefficients \eqref{eq:unit-failure} still holds.
For Gaussian coefficients we have, for all large $j$,
\begin{align*}
 \Prob(E_{K,j})
 &\le4\Pi_K(N)+5\Pi_2(N)+N^{-3}+2N^{-10}+C(K)N^{-3/8}(\log N)^4\\
 &\quad+4e^{-3N/(32e^{6K})}+5e^{-3N/(32e^{12})}+e^{-3(N+1)/32}
 +2(2N+1)e^{-N^2/2}.
\end{align*}
Similarly, for bounded coefficients we have, for all large $j$,
\begin{align*}
 \Prob(E_{K,j})
 &\le4\Pi_K(N)+5\Pi_2(N)+N^{-3}+2N^{-10}+C(K)N^{-3/8}(\log N)^4\\
 &\quad+4e^{-N/(2B^4e^{6K})}+5e^{-N/(2B^4e^{12})}+e^{-(N+1)/(2B^4)}
 +N^{-D|\log\Prob(\xi_0=0)|},
\end{align*}
where the last term is present only if $\Prob(\xi_0=0)>0$.
In both bounds, the first two exponential terms are the energy
exceptions of Proposition~\ref{prop:coefficient-transfer} at the widths
$K$ and $2$, and the third is the additional exception of
Lemma~\ref{lem:local-counts-gauss} or~\ref{lem:local-counts-bounded}.
At $N=j^8$ the terms of \eqref{eq:unit-failure} are summable as in
Step~5 of Section~\ref{sec:unit-circle}, the exponential terms by
\eqref{eq:model-summable}, and the trailing-run term by the sum above.
Hence the Borel--Cantelli argument of Step~5 applies, with the additional
term $\lceil D\log N\rceil/N\to0$ in the final bound when
$\Prob(\xi_0=0)>0$.
\end{proof}

\begin{proof}[Proof of Theorem~\ref{thm:radial-profile} for the other distributions]
We follow the proof in Section~\ref{sec:moving-radii} for Rademacher
coefficients step by step, with the normalization
$\mathbb E|\xi_0|^2=1$ and the changes written out in the proof of
Corollary~\ref{thm:stronglaw} above.
Let $\Pi_K(N)$ be as in that proof, and let $\Pi_R(N)$ be
$\Pi_K(N)$ with $R$ in place of $K$.
Fix $K,R,x$.
We use the logarithmic, mesh, local-count and tail estimates from
Section~\ref{sec:layer2-models}, and for Gaussian and bounded distributions the
probabilities of the additional exceptional events are exponential in $N$.
As in that proof, the variance profile is unchanged.
The concentration bound of Proposition~\ref{prop:coefficient-transfer}
at the five radii $r_q$ of width $R$ shows that \eqref{eq:moving-probes},
\eqref{eq:moving-secant} and \eqref{eq:moving-band} hold with the constant
$C(R)$ of the distribution, outside five events of probabilities at most
$\Pi_R(N)+\Prob(E^c)$, with $\Prob(E^c)$ at the width $R$.
The selected domains are the same as in Steps~1 and~2 of that proof.
For bounded coefficients, the band condition involves $2Ba/d_0$ instead of
$2a/d_0$, and, as in Section~\ref{sec:moving-radii}, we have
\begin{align*}
 \frac Nh\left(\frac{2Ba}{d_0}+\frac{|x|m_j}{N^2}\right)
 &\le90Be^{2(K+2)}N^{-67/64+65/64}\sqrt{\log N}+9|x|N^{7/8-2+65/64}\\
 &=90Be^{2(K+2)}N^{-1/32}\sqrt{\log N}+9|x|N^{-7/64},
\end{align*}
which tends to zero.

If $\Prob(\xi_0=0)\in(0,1)$, we take $D$ and the trailing-run event as in
that proof.
Outside that event, the degree correction $\max(0,\deg Q-\deg P)$ in
Corollary~\ref{cor:T4-two-radii} is bounded by
$m_j+\lceil D\log N\rceil$, and the extra term divided by $N$ tends to
zero.
Indeed, the three-term bound of Section~\ref{sec:moving-radii} becomes
\begin{align*}
 \left|\frac{\nu_n(1+x/n)}n-\Phi(x)\right|
 &\le\varepsilon_N
  +\frac{2\varepsilon_NN+U_N+m_j+\lceil D\log N\rceil}N+\frac{m_j}N\\
 &=3\varepsilon_N+\frac{U_N}N+\frac{2m_j}N+\frac{\lceil D\log N\rceil}N.
\end{align*}
For Gaussian coefficients we have, for all large $j$,
\begin{align*}
 \Prob(E_{K,R,x,j})
 &\le4\Pi_K(N)+5\Pi_R(N)+N^{-3}+2N^{-10}+C(K)N^{-3/8}(\log N)^4\\
 &\quad+4e^{-3N/(32e^{6K})}+5e^{-3N/(32e^{6R})}+e^{-3(N+1)/32}
 +2(2N+1)e^{-N^2/2},
\end{align*}
and for bounded coefficients,
\begin{align*}
 \Prob(E_{K,R,x,j})
 &\le4\Pi_K(N)+5\Pi_R(N)+N^{-3}+2N^{-10}+C(K)N^{-3/8}(\log N)^4\\
 &\quad+4e^{-N/(2B^4e^{6K})}+5e^{-N/(2B^4e^{6R})}+e^{-(N+1)/(2B^4)}
 +N^{-D|\log\Prob(\xi_0=0)|},
\end{align*}
where the last term is present only if $\Prob(\xi_0=0)>0$.
These bounds are obtained from those of Step~5 in that proof by
replacing the five thin-band terms at the width $2$ with the five logarithmic terms
at the width $R$, as \eqref{eq:moving-failure} replaces
\eqref{eq:unit-failure}.
For Steinhaus coefficients \eqref{eq:moving-failure} still holds.
As in that proof, these bounds are summable along $N=j^8$.
The conclusion now follows from the fixed-$K$ intersection and the
rational-$x$ argument of Section~\ref{sec:moving-radii}.
\end{proof}
